\begin{figure}[!htbp]
\centering
\begin{tikzpicture}
\begin{axis}[
  width=1.02\columnwidth, height=46mm,
  ybar, bar width=3.4pt,
  enlarge x limits=0.28,
  ymin=0.9, ymax=2.05,
  ytick={1.0,1.2,1.4,1.6,1.8,2.0},
  symbolic x coords={1,4,16},
  xtick=data,
  xlabel={\ifSpanish Filas $N$\else Rows $N$\fi},
  ylabel={\ifSpanish Ciclos B3 / ciclos D\else B3 cycles / D cycles\fi},
  tick label style={font=\scriptsize},
  label style={font=\scriptsize},
  legend style={font=\scriptsize,at={(0.5,1.03)},anchor=south,
                legend columns=4,draw=none,column sep=3pt},
  ymajorgrids, grid style={gray!25},
  axis line style={black}, tick style={black},
]
\addplot[fill=white,draw=black] coordinates {(1,1.00) (4,1.00) (16,1.00)};
\addplot[fill=gray!35,draw=black] coordinates {(1,1.86) (4,1.25) (16,1.08)};
\addplot[fill=gray!70,draw=black] coordinates {(1,1.88) (4,1.26) (16,1.08)};
\addplot[fill=black,draw=black] coordinates {(1,1.00) (4,1.46) (16,1.28)};
\legend{$z{=}0$, $z{=}8$, $z{=}15$, ZS}
\draw[dashed,gray] ({rel axis cs:0,0}|-{axis cs:1,1.0})
                -- ({rel axis cs:1,0}|-{axis cs:1,1.0});
\end{axis}
\end{tikzpicture}
\caption{\ifSpanish Ventaja de la fusión por régimen y número de filas. La línea
punteada en 1.0 marca la ausencia de ventaja. Los tres primeros grupos son
puntos cero constantes; ZS los hace variar por fila. La ventaja se estrecha al
crecer $N$ porque B3 amortiza su corrección entre más filas, y en $z{=}0$ no
existe en absoluto porque B3 la omite. En $N{=}1$ un punto cero variable es
degenerado: una sola fila no puede tener un $z$ que varíe. Corrida
\CampaignRun.
\else Advantage of fusion by regime and row count. The dashed line at 1.0 marks
the absence of any advantage. The first three groups are constant zero points;
ZS makes them vary per row. The advantage narrows as $N$ grows because B3
amortizes its correction over more rows, and at $z{=}0$ there is none at all
because B3 omits it. At $N{=}1$ a varying zero point is degenerate: a single row
cannot have a $z$ that varies. Run \CampaignRun.\fi}
\label{fig:ratio}
\end{figure}
